\section{Implementation}
\label{sec:implementation}




In \sysname{} we use Qwen2.5-VL-7B for constructing EKGs, Qwen2.5-32B for SA, and Gemini-1.5-Pro for CA. \revised{We utilize AWQ~\cite{lin2024awq} and LMDeploy~\cite{2023lmdeploy} to accelerate the on-device inference. These particular LLMs are selected for \sysname due to our two-step principle: we first determine an appropriate model size that is feasible for edge deployment, and then, within that size range, choose the best-performing model available according to the public benchmarks~\cite{fu2024videomme, shen2024longvu}.} Additionally, \sysname{} adopts batch inference for several key stages—including description generation, description merging, entity extraction, and tree search—to improve efficiency and maximize GPU utilization. For text and vision embedding, we utilize JinaCLIP\cite{koukounas2024jinaclip}.
And we employ BERTScore with the deberta-xlarge-mnli\cite{he2020deberta} checkpoint. The storage of EKG and vector representations is based on the implementation of \cite{guo2024lightrag}, upon which we make further modifications to suit the specific requirements of \sysname{}. 

During the EKGs construction stage, we carefully design prompts to guide the extraction of structured information from video. For general-purpose video understanding, we employ a unified prompt that avoids introducing bias or prior assumptions:

\textit{"...Your task is to extract and provide a detailed description of the video segment, focusing on all key visible details..."}. 

For scenario-specific videos, we design the prompts to emphasize scenario-relevant information. \revised{For example, in the case of wildlife Monitoring scenario, key information may includes the timestamps of the recording, animal activities (\eg presence, species, number, specific behaviors, etc.), and environmental changes. Particularly the prompts used in \sysname are listed in \S\ref{app:prompts}}

\revised{The hyperparameters in \sysname are determined through a parameter tuning procedure conducted on the selected datasets. Specifically, for the semantic chunking methodology described in \S\ref{subsec: sematic_chunking}, we set the BERTScore threshold for merging and summarizing descriptions at 0.65. Empirically, this threshold yielded partitions that exhibited greater alignment with human annotations across various scenarios. Additionally, the maximum tree search depth is set to 3, more details are provided in \S\ref{subsub:depth}.}




